% optimization/alignment/50_sqp_analytic_structure.tex

\section{Analytical Structure of the Optimization Problem}

% ============================================================
\subsection{Motivation}
% ============================================================

In classical numerical optimization, Jacobians are often computed by
finite differences or automatic differentiation.

Within the AIM framework, this is not necessary for the reconstruction
block: the planar reconstruction is differentiable in its parameters
under explicit hypotheses stated below.

% ============================================================
\subsection{Fundamental Observation}
% ============================================================

\begin{corollary}[Planar reconstruction integrals]
\label{cor:planar-reconstruction-integrals}
By \cref{prop:curvature-driven-reconstruction}, the planar pose2 variables
are obtained from the curvature law by arc-length integration:
\[
\theta(s) = \theta_0 + \int_0^s \kappa(\sigma)\,\dd\sigma,
\]
\[
\gamma(s) = \gamma_0 + \int_0^s (\cos\theta,\sin\theta)\,\dd\sigma.
\]
\end{corollary}

Profile, cant, chainage, and realization state are not obtained from these
two integrals; they require their own longitudinal laws and operators.

% ============================================================
\subsection{Parameter Types}
% ============================================================

\begin{definition}
Primary parameter classes:
\begin{itemize}
\item curvature parameters \(dK\),
\item length parameters \(dL\),
\item transition parameters \(dT\).
\end{itemize}
\end{definition}

% ============================================================
\subsection{Analytical Derivatives}
% ============================================================

\begin{proposition}[Parameter derivative of heading]
\label{prop:heading-parameter-derivative}
Let \(\kappa(\sigma,p)\) be differentiable in the parameter \(p\) for
almost every \(\sigma\), and let
\(\lvert\partial_p\kappa(\sigma,p)\rvert\le m(\sigma)\) locally uniformly
in \(p\) with \(m\in L^1(0,s)\). Then
\[
\frac{\partial \theta(s)}{\partial p}
=
\int_0^s
\frac{\partial \kappa}{\partial p}\,\dd\sigma.
\]
\end{proposition}

\begin{proof}
Under the stated domination hypothesis, differentiation under the integral
sign applies to
\(\theta(s,p)=\theta_0+\int_0^s\kappa(\sigma,p)\,\dd\sigma\).
\end{proof}

\begin{proposition}[Parameter derivative of position]
\label{prop:position-parameter-derivative}
Under the hypotheses of \cref{prop:heading-parameter-derivative},
\[
\frac{\partial \gamma(s)}{\partial p}
=
\int_0^s
\begin{pmatrix}
-\sin\theta \\
\cos\theta
\end{pmatrix}
\frac{\partial \theta}{\partial p}\,\dd\sigma.
\]
\end{proposition}

\begin{proof}
The chain rule applied to
\(\sigma\mapsto(\cos\theta(\sigma,p),\sin\theta(\sigma,p))\) gives the
integrand \((-\sin\theta,\cos\theta)^{\!\top}\partial_p\theta\). By
\cref{prop:heading-parameter-derivative},
\(\lvert\partial_p\theta\rvert\le\lVert m\rVert_{L^1(0,s)}\) on \([0,s]\),
so the integrand is dominated and differentiation under the integral sign
is admissible.
\end{proof}

Within these hypotheses, all derivatives reduce to arc-length integrals.

% ============================================================
\subsection{Key Insight}
% ============================================================

\begin{proposition}[Exact Jacobians of the smooth reconstruction block]
\label{prop:exact-jacobians-smooth-block}
For residuals and constraints that are differentiable compositions of the
parameterization, \(\theta\), and \(\gamma\), the Jacobian is given exactly
by
\cref{prop:heading-parameter-derivative,prop:position-parameter-derivative}
and requires no finite-difference approximation.
\end{proposition}

The statement is restricted to the smooth reconstruction block.
Nondifferentiable admission logic, active-set changes of inequality
constraints, and black-box operators do not automatically admit exact
Jacobians and must be treated separately.

Within that block, this removes:
\begin{itemize}
\item finite differences and their truncation error,
\item repeated reconstruction for each perturbed parameter,
\item the associated computational cost.
\end{itemize}

% ============================================================
\subsection{Block Structure}
% ============================================================

The parameter vector decomposes as:
\[
(\boldsymbol{\kappa}, \boldsymbol{\phi}).
\]

This induces a block structure:
\[
H =
\begin{pmatrix}
H_{\kappa\kappa} & H_{\kappa\phi} \\
H_{\phi\kappa} & H_{\phi\phi}
\end{pmatrix}.
\]

% ============================================================
\subsection{Interpretation}
% ============================================================

\begin{summary}
The structure of the optimization problem is not imposed numerically,
but follows directly from the curvature-based modelling approach.

This yields:
\begin{itemize}
\item exact derivatives,
\item sparse matrices,
\item natural block structure.
\end{itemize}

This is a central advantage of the AIM framework.
\end{summary}
